\scanpage{300}
\mspage{41}
\[
\begin{aligned}
\mathrm{Tr}^+(f)^*&\colon\mathrm{Tr}^+(\unsure{\tau|t}_1)\unsure{\mathbin{-}|\longrightarrow}\mathrm{Tr}^+(\tau_2)\\
\mathrm{Tr}^-(f)^*&\quad\mathrm{Tr}^{\deleted{\unsure{+|-}}\unsure{-|+}}(\tau_1)\mathbin{-}\mathrm{Tr}^-(\tau_{\unsure{1|2}}),
\end{aligned}
\]
en d'autres termes, une application
\[
 \mathrm{Tr}(f)^{\unsure{*|'}}\colon\mathrm{Tr}(\tau_1)\mathbin{-}\mathrm{Tr}(\tau_2)
\]
``compatible avec les signatures''. Dans le cor. 5
plus haut, on a défini une telle application seulement pour $\tau'$ réduit.
% slanted
\margin{Ce qui suit ne fait pas intervenir les signatures -- c'est un \unsure{sorite|sorites} sur les applications croissantes entre \unsure{ens.} finis totalement ordonnés.}
Pour une flèche \deleted{\unsure{de $f$}}
de $I_{\unsure{\tau,|\tau_1}}$, \deleted{désign} appelons \uline{portée ensembliste} de
$f$ sur $u$, ou de $u$ par $f$, la portée \inserted{ensembliste} de
la flèche $f(u)$ dans $I_{\tau'}$, i.e. \deleted{\unsure{le segment}}
le segment $[f(i-1),f(i)]\subset I_{\tau'}$ (si $u=\underline{u}_{\unsure{i}}$).
\deleted{On voit alors que les portées des}

Soient $\underline{u}_{i_1},\ldots,\underline{u}_{i_N}$ ou $(v_1,v_2,\ldots,v_N)$
la suite croissante des éléments de $\mathrm{Tr}(\tau,f)$,
on voit \inserted{\unsure{donc par réc. sur $i$|disons par réc. sur $i$}} que pour \deleted{\unsure{$1<i\leq N$}} $i\in[1,N-1]$, on a
\[
\left\{\begin{aligned}
 f(\operatorname{but}(v_i))&=f(\operatorname{source}(v_{i+1}))\deleted{\unsure{=}}\ ( = t_i)\\
 \deleted{\text{et si on désigne}}\quad f(\operatorname{source}v_1)&=0_{\tau'},\quad f(\operatorname{but}(v_N))=1_{\tau'}^{\unsure{+}},
\end{aligned}\right.
\]
Considérons la suite
\[
\{\overset{\substack{t_0\\\shortparallel}}{0_{\tau'}},t_1,t_2,\ldots\ t_{N-1},\overset{\substack{t_N\\\shortparallel}}{1_{\tau'}}\}
\]
elle est définie par le sous-ensemble correspondant de $I_{\tau'}$, qu'on appellera l'ensemble des
\uline{sommets critiques} \uline{de} $I_{\tau'}$ \uline{pour} $f$. \deleted{Donc} C'est
une suite \inserted{contenant $0_{\tau'},1_{\tau'}$ et} telle que deux éléments consécutifs
sont reliés dans $I_{\tau'}$ par une flèche, ce qui \deleted{implique}
\inserted{équivaut à dire} qu'elle contient les sommets critiques (intrinsèques)
de $I_{\tau'}$. \deleted{\unsure{Or}} D'autre part, pour les \deleted{flèches de}\inserted{intervalles}
de transition $[i-1,i]$ de $I_{\tau'}$ \deleted{(qui} $(i\in[1,n'])$,
% archive sheet number: 299
